% Claim statement file for row claim_3_13 -- "AcyclicHardInterventionTopologicalOrder".
% Auto-generated stub by scaffold/solve_chapter.py at row start. The
% manager agent (or a worker it dispatches) fills the body below with the
% claim's statement(s) from the lecture notes -- statement only, no proof
% (proofs live in the sibling `claim_3_13_proof_*.tex` file).
%
% This file is a sibling of `main.tex` inside `<subsection>/tex/`. The
% `subfiles` package lets it render both standalone and as part of
% `main.tex`. Note: `main.tex` does NOT \subfile this statement file -- the
% sibling `_proof_` file restates the statement above the proof, so
% including both in `main.tex` would be redundant. This file exists for
% standalone reading / grep convenience.

\documentclass[main]{subfiles}

\begin{document}
% ref: claim_3_13 (statement)
% title: AcyclicHardInterventionTopologicalOrder
\def\rowref{claim\_3\_13}
\phantomsection\label{claim_3_13}

\begin{Rem}[AcyclicHardInterventionTopologicalOrder]
    Let $G = (J, V, E, L)$ be a CDMG (def \ref{def-cdmg}); throughout we use the notation of def \ref{not-cdmg} and recall the typing and axioms of def \ref{def-cdmg}, namely $J \cap V = \emptyset$, $E \ins (J \cup V) \times V$, $L \ins V \times V$, $L$ irreflexive ($(v_1, v_2) \in L \Rightarrow v_1 \neq v_2$), and $L$ symmetric ($(v_1, v_2) \in L \Leftrightarrow (v_2, v_1) \in L$). Let
    \[ W \ins J \cup V \]
    be a subset of nodes, and let
    \[ G_{\doit(I_W)} \;=\; \lp J_{\doit(I_W)},\, V_{\doit(I_W)},\, E_{\doit(I_W)},\, L_{\doit(I_W)} \rp \]
    be the extended CDMG of $G$ w.r.t.\ nodes $W$ and corresponding intervention nodes $I_W$, as constructed in def \ref{def:cdmg_intervention_nodes}. Both quantifiers ``for every CDMG $G$'' and ``for every $W \ins J \cup V$'' are read \emph{universally}, with the side condition $W \ins J \cup V$ inherited verbatim from def \ref{def:cdmg_intervention_nodes} (so $W$ is permitted to overlap the input-node set $J$, i.e.\ $W \cap J \neq \emptyset$ is admitted); the case $W = \emptyset$ is admitted as well (then $\lC I_w \st w \in W \sm J \rC = \emptyset$ and $G_{\doit(I_W)} = G$ verbatim by items i.--iv.\ of def \ref{def:cdmg_intervention_nodes}).

    \emph{Notational shorthand recap from def \ref{def:cdmg_intervention_nodes}.}
    Per def \ref{def:cdmg_intervention_nodes}, the family of intervention symbols is $I_W = \lC I_w \st w \in W \rC$ with the convention $I_j := j$ for $j \in J \cap W$ (a purely notational shorthand: no new node is introduced on the $J \cap W$ branch, the pre-existing input node $j \in J$ plays the role of its own intervention node, and the pre-existing edges incident to $j$ in $G$ are retained unchanged in $G_{\doit(I_W)}$). Consequently $I_W$ decomposes as
    \[
        I_W \;=\; (J \cap W) \;\cup\; \lC I_w \st w \in W \sm J \rC,
    \]
    where the first component contains the (notational) $I_j = j$ for $j \in J \cap W$ and the second component contains the genuinely fresh symbols $I_w$ for $w \in W \sm J$, the latter type-disjoint from $J \cup V$ and from each other by the freshness clause of def \ref{def:cdmg_intervention_nodes}. We henceforth refer to the elements of $\lC I_w \st w \in W \sm J \rC$ as the \emph{fresh intervention nodes} of $G_{\doit(I_W)}$.

    \emph{Carrier of $G_{\doit(I_W)}$.}
    By items i.\ and ii.\ of def \ref{def:cdmg_intervention_nodes},
    \[
        J_{\doit(I_W)} \cup V_{\doit(I_W)} \;=\; \lp J \cup \lC I_w \st w \in W \sm J \rC \rp \cup V \;=\; (J \cup V) \;\dcup\; \lC I_w \st w \in W \sm J \rC,
    \]
    where the second equality is a type-level disjoint union by the freshness clause of def \ref{def:cdmg_intervention_nodes} ($\lC I_w \st w \in W \sm J \rC \cap (J \cup V) = \emptyset$). In particular, $J \cup V$ is a subset of $J_{\doit(I_W)} \cup V_{\doit(I_W)}$ via the canonical set-theoretic inclusion
    \[
        \iota \,:\, J \cup V \;\hookrightarrow\; J_{\doit(I_W)} \cup V_{\doit(I_W)}, \qquad v \;\longmapsto\; v,
    \]
    and the inclusion is strict when $W \sm J \neq \emptyset$ (the fresh intervention nodes belong to the codomain but not to the domain) and is an equality when $W \sm J = \emptyset$. In the downstream Lean rendering of def \ref{def:cdmg_intervention_nodes} this inclusion is realised as the ``unsplit'' constructor of a tagged-sum carrier; the underlying set-theoretic content is plain inclusion.

    \emph{Topological orders live on different carriers.}
    By def \ref{def-topological-order}, a topological order of $G$ is a strict total order $<_G$ on $J \cup V$, while a topological order of $G_{\doit(I_W)}$ is a strict total order $<_{G_{\doit(I_W)}}$ on $J_{\doit(I_W)} \cup V_{\doit(I_W)}$. When $W \sm J \neq \emptyset$, these carriers are not equal -- $<_{G_{\doit(I_W)}}$ is a relation on the strictly larger set $J_{\doit(I_W)} \cup V_{\doit(I_W)}$ -- so a topological order of $G_{\doit(I_W)}$ is not literally a topological order of $G$ in the sense of def \ref{def-topological-order}; the bridging operations are \emph{restriction} (sub-claim~(b) below) and \emph{extension} (sub-claim~(c) below), both taken along $\iota$.

    \emph{Reading of ``extend''.}
    Following standard mathematical usage, we read ``extension of a strict total order'' in the \emph{strict} sense: a strict total order $\prec$ on $J_{\doit(I_W)} \cup V_{\doit(I_W)}$ \emph{extends} a strict total order $<_G$ on $J \cup V$ when its restriction to $J \cup V$ along $\iota$ recovers $<_G$, i.e.\ for every $v_1, v_2 \in J \cup V$,
    \[
        v_1 \,<_G\, v_2 \;\Longleftrightarrow\; \iota(v_1) \,\prec\, \iota(v_2).
    \]
    This is the strict reading; the weaker reading ``some strict total order on $J_{\doit(I_W)} \cup V_{\doit(I_W)}$ whose ground set contains $J \cup V$'' is too lax to constrain the construction and is not adopted here.

    \emph{The remark.} The following conjunction of three sub-claims holds.

    \begin{enumerate}[label=(\alph*)]
        \item \emph{Acyclicity is preserved by adding intervention nodes.}
              If $G$ is acyclic in the sense of def \ref{def-acylic}, then $G_{\doit(I_W)}$ is acyclic in the sense of def \ref{def-acylic}, i.e.\ for every $x \in J_{\doit(I_W)} \cup V_{\doit(I_W)} = (J \cup V) \dcup \lC I_w \st w \in W \sm J \rC$, there does not exist any non-trivial directed walk from $x$ to itself in $G_{\doit(I_W)}$. The hypothesis ``$G$ is acyclic'' is used in this sub-claim; sub-claims~(b) and~(c) below do not require acyclicity of $G$.

        \item \emph{Restriction direction: every topological order of $G_{\doit(I_W)}$ restricts to one of $G$ along $\iota$.}
              For every strict total order $\prec$ on $J_{\doit(I_W)} \cup V_{\doit(I_W)}$, if $\prec$ is a topological order of $G_{\doit(I_W)}$ in the sense of def \ref{def-topological-order} (i.e.\ for every $v, w \in J_{\doit(I_W)} \cup V_{\doit(I_W)}$, $v \in \Pa^{G_{\doit(I_W)}}(w) \Rightarrow v \prec w$), then the binary relation $<_G$ on $J \cup V$ defined by
              \[
                  v_1 \,<_G\, v_2 \;\;:\Longleftrightarrow\;\; \iota(v_1) \,\prec\, \iota(v_2) \qquad \text{for } v_1, v_2 \in J \cup V
              \]
              is a topological order of $G$ in the sense of def \ref{def-topological-order}. The LN's prose ``a topological order for $G_{\doit(I_W)}$ is also one for $G$'' is read \emph{universally} (every such $\prec$, not merely some such $\prec$) and refers to this restriction along $\iota$: the carriers differ when $W \sm J \neq \emptyset$, so an order on $J_{\doit(I_W)} \cup V_{\doit(I_W)}$ cannot literally be an order on $J \cup V$, but its restriction along $\iota$ is.

        \item \emph{Extension direction: every topological order of $G$ extends to one of $G_{\doit(I_W)}$.}
              For every strict total order $<_G$ on $J \cup V$, if $<_G$ is a topological order of $G$ in the sense of def \ref{def-topological-order}, then there exists a strict total order $\prec$ on $J_{\doit(I_W)} \cup V_{\doit(I_W)}$ such that
              \begin{itemize}
                  \item $\prec$ is a topological order of $G_{\doit(I_W)}$ in the sense of def \ref{def-topological-order}; and
                  \item $\prec$ \emph{extends} $<_G$ along $\iota$ in the strict sense recorded above, i.e.\ for every $v_1, v_2 \in J \cup V$, $v_1 <_G v_2 \Leftrightarrow \iota(v_1) \prec \iota(v_2)$.
              \end{itemize}
              The LN's prose ``any topological order of $G$ can be extended to one for $G_{\doit(I_W)}$'' is read \emph{universally} (every such $<_G$ extends) and ``extended'' is read in the strict sense.

              \emph{Example construction (the LN's ``put all $I_w$ first'').}
              A concrete witness $\prec$ for the existential of sub-claim~(c) is the following ``fresh intervention nodes first'' relation: for $x, y \in J_{\doit(I_W)} \cup V_{\doit(I_W)} = (J \cup V) \dcup \lC I_w \st w \in W \sm J \rC$,
              \[
                  x \,\prec\, y \;\;:\Longleftrightarrow\;\;
                  \begin{cases}
                      w_1 \,<_G\, w_2 & \text{if } x = I_{w_1} \text{ and } y = I_{w_2} \text{ for some } w_1, w_2 \in W \sm J, \\
                      \text{true} & \text{if } x = I_{w} \text{ for some } w \in W \sm J \text{ and } y \in J \cup V, \\
                      \text{false} & \text{if } x \in J \cup V \text{ and } y = I_{w} \text{ for some } w \in W \sm J, \\
                      v_1 \,<_G\, v_2 & \text{if } x = v_1 \text{ and } y = v_2 \text{ for some } v_1, v_2 \in J \cup V.
                  \end{cases}
              \]
              The four cases exhaust the disjoint-union carrier $(J \cup V) \dcup \lC I_w \st w \in W \sm J \rC$ and are pairwise mutually exclusive, so the right-hand side is well-defined. In the first case the right-hand side $w_1 <_G w_2$ uses that $W \sm J \ins V \ins J \cup V$ -- both target nodes $w_1, w_2$ lie in $V \ins J \cup V$, the carrier of $<_G$. In words: all fresh intervention nodes $I_w$ (i.e.\ those with $w \in W \sm J$) precede all elements of $J \cup V$ under $\prec$; among the fresh intervention nodes, $\prec$ is the $<_G$-order of their target nodes; among the elements of $J \cup V$, $\prec$ is $<_G$ itself. The restriction of $\prec$ to $J \cup V$ along $\iota$ is the fourth case, which is $<_G$ verbatim, so the strict extension property is satisfied by construction.
    \end{enumerate}

    \emph{What ``putting all the $I_w$ nodes first'' refers to (resolving the literal reading of the LN's prose).}
    By the notational shorthand $I_j := j$ for $j \in J \cap W$ recalled above from def \ref{def:cdmg_intervention_nodes}, the literal LN phrase ``putting all the $I_w$ nodes first'' must be read as ``putting all the \emph{fresh} intervention nodes first'', i.e.\ only the $I_w$ with $w \in W \sm J$ are placed at the front of the ordering; the $I_j = j$ for $j \in J \cap W$ are not separate nodes, they coincide with pre-existing input nodes of $G$, and they retain their $<_G$-position inside the $J \cup V$-block of $\prec$. This is the only reading consistent with the strict ``extends'' property of the previous paragraph: were the $J \cap W$ nodes also moved to the front, then the restriction of $\prec$ to $J \cup V$ along $\iota$ would in general differ from $<_G$ (it would reshuffle the $J \cap W$ nodes), violating the strict extension condition.

    \emph{Order among the fresh intervention nodes is a design freedom.}
    The LN's prose does not specify the relative order \emph{among} the fresh intervention nodes $\lC I_w \st w \in W \sm J \rC$. Any choice that makes $\prec$ a strict total order on $J_{\doit(I_W)} \cup V_{\doit(I_W)}$ and respects the parent-precedence requirement of def \ref{def-topological-order} for $G_{\doit(I_W)}$ is admissible. The example construction above selects the specific choice ``$I_{w_1} \prec I_{w_2} \Leftrightarrow w_1 <_G w_2$'' (using that $W \sm J \ins V \ins J \cup V$, so $<_G$ is defined on the targets); but any strict total order on $\lC I_w \st w \in W \sm J \rC$ would equally well discharge the topological-order requirement for $G_{\doit(I_W)}$, because the fresh intervention nodes are pairwise non-adjacent in $G_{\doit(I_W)}$: each fresh $I_w$ has the unique outgoing edge $(I_w, w) \in E_{\doit(I_W)}$ pointing to its target $w \in V$, has no incoming edges in $G_{\doit(I_W)}$ (by the freshness clause of def \ref{def:cdmg_intervention_nodes}, no element of $E_{\doit(I_W)}$ has a fresh $I_w$ as its target), and the set $\lC (I_w, w) \st w \in W \sm J \rC$ contains no pair $(I_{w_1}, I_{w_2})$ with both endpoints fresh. Sub-claim~(c) asserts \emph{existence} of an extension; the example construction above provides one explicit witness.

    \emph{Role of $L$.}
    The bidirected-edge set $L$ plays no role in any of the three sub-claims: acyclicity (def \ref{def-acylic}) is defined via the non-existence of non-trivial directed walks (which use only $E$), and a topological order (def \ref{def-topological-order}) is defined via the parent-set $\Pa^G$ (def \ref{def_3_5}, item~i., which is defined from $E$ alone). By item iv.\ of def \ref{def:cdmg_intervention_nodes}, $L_{\doit(I_W)} = L$ is transported verbatim, but neither sub-claim consults it.
\end{Rem}

\end{document}
